\chapter{Complexity — Encoding}
%%% generated by the blueprint pipeline from TCSlib.Complexity.TuringMachine.Encoding
%%%%%%%%%%%%%%%%%%%%%%%%%%%%%%%%%%%%%%%%%%%%%%%%%%%%%%%%%%%%%%%%%%%%%%%%%%%%%%%
\section{Overview}

This module represents Turing machines as binary strings (Arora--Barak,
\S 1.4): it defines the \emph{code normal form} (one work tape, binary
alphabet, states drawn from a canonical nonempty finite type), a fixed
canonical serialization of such machines, the self-delimiting pair encoding
used by the universal machine together with a concrete linear-time machine
computing the diagonal pairing, and the specification of a machine
representation scheme --- both its algebraic laws (every string decodes to a
machine, and codes survive padding) and the effectivity requirement that ties
decoding to computable semantics.  Throughout, bit strings are written over
$\{0,1\}$ ($1$ for true, $0$ for false), and the \emph{pair encoding}
$\langle x, \alpha\rangle$ of two bit strings lists each bit of $x$ twice,
then the unequal separator pair $01$, then $\alpha$ verbatim.

%%%%%%%%%%%%%%%%%%%%%%%%%%%%%%%%%%%%%%%%%%%%%%%%%%%%%%%%%%%%%%%%%%%%%%%%%%%%%%%
\section{Declarations}

\begin{definition}[Machine in code normal form]
\lean{Turing.CodeTM}
\label{Turing.CodeTM}
\leanok
\uses{Turing.MultiTapeTM}
A machine in \emph{code normal form} consists of a natural number $m$ (one
less than the number of states, so the state space is never empty) together
with a multi-tape Turing machine having exactly one work tape, binary tape
alphabet $\{0,1\}$, and state space the canonical finite type
$\{0, 1, \dots, m\}$.
\end{definition}

\begin{definition}[Bundled machine of a coded machine]
\lean{Turing.CodeTM.toFinTM}
\label{Turing.CodeTM.toFinTM}
\leanok
\uses{Turing.CodeTM, Turing.FinTM}
The finite-machine bundle associated with a machine in code normal form: the
binary-alphabet machine package whose number of work tapes is $1$, whose
state type is the canonical finite state space $\{0,\dots,m\}$ of the coded
machine, and whose transition table is that of the coded machine.
\end{definition}

\begin{lemma}[Tape count of the bundled machine]
\lean{Turing.CodeTM.toFinTM_k}
\label{Turing.CodeTM.toFinTM_k}
\leanok
\uses{Turing.CodeTM, Turing.CodeTM.toFinTM}
\difficulty{1}
For every machine $M$ in code normal form, the bundled finite machine
associated with $M$ has exactly one work tape.
\end{lemma}

\begin{definition}[Self-delimiting pair encoding]
\lean{Turing.pairEncode}
\label{Turing.pairEncode}
\leanok
The self-delimiting pairing $\langle x, \alpha\rangle$ of two bit strings:
the first string $x$ with every bit doubled, then the separator pair $01$,
then the second string $\alpha$ verbatim.  Reading aligned two-bit blocks
from the left, $00$ and $11$ are data and the first aligned $01$ is the
separator (an aligned $01$ can never occur inside the doubled region), so the
encoding determines both components.
\end{definition}

\begin{definition}[Aligned pair parser]
\lean{Turing.pairDecode}
\label{Turing.pairDecode}
\leanok
A partial parser on bit strings that reads aligned two-bit blocks from the
left: a block $00$ contributes the bit $0$ and a block $11$ contributes the
bit $1$ to the first component, and the first aligned block $01$ is the
separator, at which point the parser returns the accumulated first component
together with the remaining suffix, untouched, as the second component.  On
strings admitting no such parse the result is undefined.
\end{definition}

\begin{lemma}[The parser inverts the pair encoding]
\lean{Turing.pairDecode_pairEncode}
\label{Turing.pairDecode_pairEncode}
\leanok
\uses{Turing.pairDecode, Turing.pairEncode}
\difficulty{4}
For all bit strings $x$ and $\alpha$, the aligned two-bit parser applied to
the pair encoding $\langle x, \alpha\rangle$ (each bit of $x$ doubled, then
the separator $01$, then $\alpha$ verbatim) succeeds and returns exactly the
pair $(x, \alpha)$.
\end{lemma}

\begin{theorem}[Injectivity of the pair encoding]
\lean{Turing.pairEncode_injective}
\label{Turing.pairEncode_injective}
\leanok
\uses{Turing.pairDecode, Turing.pairDecode_pairEncode, Turing.pairEncode}
\difficulty{3}
The map sending a pair $(x, \alpha)$ of bit strings to its pair encoding ---
each bit of $x$ doubled, then the separator $01$, then $\alpha$ verbatim ---
is injective on pairs of bit strings.
\end{theorem}

\begin{definition}[The diagonal pairing machine]
\lean{Turing.pairDiagTM}
\label{Turing.pairDiagTM}
\leanok
\uses{Turing.FinTM}
A fixed six-state Turing machine over the binary alphabet with no work tapes
that transforms an input $x$ into the diagonal pair encoding
$\langle x, x\rangle$.  Its six control states implement, in order: emitting
the scanned bit while staying put, emitting it again while moving right,
emitting the separator bit $1$ while staying put, taking one unconditional
left move, rewinding the input head to the left end, and copying the input
verbatim; the first state's blank branch emits the separator bit $0$.
\end{definition}

\begin{definition}[Configurations of the pairing machine]
\lean{Turing.pairDiagCfg}
\label{Turing.pairDiagCfg}
\leanok
\uses{Turing.Cfg}
A configuration of the six-state diagonal pairing machine on an input string
$x$, specified by an optional control state among the six (absent when the
machine has halted), an input-head position in $\{0, \dots, |x| + 1\}$, and
the output string emitted so far.  The machine has no work tapes, so the
work-tape components of the configuration are vacuous.
\end{definition}

\begin{lemma}[One step of the pairing machine]
\lean{Turing.pairDiag_step}
\label{Turing.pairDiag_step}
\leanok
\uses{Turing.Action.apply, Turing.Cfg.ext_zero_tapes, Turing.Cfg.inputSymbol,
  Turing.Cfg.workTapeSymbols, Turing.MultiTapeTM.step, Turing.moveInputPos,
  Turing.pairDiagCfg, Turing.pairDiagTM}
\difficulty{3}
Consider the six-state diagonal pairing machine on an input string $x$, in a
live configuration with control state $q$, input-head position $p$, and
output string $w$, and suppose the symbol scanned by the input head is $b$
(a bit or the blank).  Then one step of the machine yields exactly the
configuration prescribed by the transition table entry at $(q, b)$: its
successor state (or halting), the head position $p$ moved by the prescribed
direction, and the output $w$ extended by the emitted bit, if any.
\end{lemma}

\begin{lemma}[Scanned symbol at an interior position]
\lean{Turing.pairDiag_inner}
\label{Turing.pairDiag_inner}
\leanok
\uses{Turing.Cfg.inputSymbol, Turing.inputSymbolInner, Turing.pairDiagCfg}
\difficulty{1}
For the six-state diagonal pairing machine on an input string $x$, in any
configuration whose input head sits at position $j + 1$ with $j < |x|$, the
scanned input symbol is the $j$-th bit of $x$, for every control state and
every output string.
\end{lemma}

\begin{lemma}[Scanned symbol at the right boundary]
\lean{Turing.pairDiag_right}
\label{Turing.pairDiag_right}
\leanok
\uses{Turing.Cfg.inputSymbol, Turing.pairDiagCfg}
\difficulty{1}
For the six-state diagonal pairing machine on an input string $x$, in any
configuration whose input head sits at the right boundary position
$|x| + 1$, the scanned input symbol is the blank, for every control state and
every output string --- including on the empty input.
\end{lemma}

\begin{lemma}[The doubling pass]
\lean{Turing.pairDiag_double}
\label{Turing.pairDiag_double}
\leanok
\uses{Turing.Cfg.ext_zero_tapes, Turing.MultiTapeTM.initCfg,
  Turing.MultiTapeTM.runFrom, Turing.MultiTapeTM.runFrom_succ_eq_step',
  Turing.moveInputPos_pos_of_ne_right, Turing.moveInputPos_zero,
  Turing.pairDiagCfg, Turing.pairDiagTM, Turing.pairDiag_inner,
  Turing.pairDiag_step}
\difficulty{5}
Let $x$ be a bit string and $t \le |x|$.  Running the six-state diagonal
pairing machine from its initial configuration on input $x$ for exactly $2t$
steps yields the configuration in the initial control state with input head
at position $t + 1$ whose output is the first $t$ bits of $x$ with each bit
doubled.
\end{lemma}

\begin{lemma}[The rewind pass]
\lean{Turing.pairDiag_rewind}
\label{Turing.pairDiag_rewind}
\leanok
\uses{Turing.Cfg.ext_zero_tapes, Turing.Cfg.inputSymbol,
  Turing.MultiTapeTM.runFrom, Turing.MultiTapeTM.runFrom_succ_eq_step,
  Turing.MultiTapeTM.runFrom_succ_eq_step', Turing.MultiTapeTM.runFrom_zero,
  Turing.moveInputPos_neg_of_ne_left, Turing.moveInputPos_pos_of_ne_right,
  Turing.pairDiagCfg, Turing.pairDiagTM, Turing.pairDiag_inner,
  Turing.pairDiag_step}
\difficulty{4}
Let $x$ be a bit string, $w$ an output string, and $j \le |x|$.  Running the
six-state diagonal pairing machine for exactly $j + 1$ steps from the
configuration in its rewind state (state $4$) with input head at position $j$
and output $w$ yields the configuration in its copy state (state $5$) with
input head at position $1$ and output still $w$.
\end{lemma}

\begin{lemma}[The copy pass]
\lean{Turing.pairDiag_copy}
\label{Turing.pairDiag_copy}
\leanok
\uses{Turing.Cfg.ext_zero_tapes, Turing.MultiTapeTM.runFrom,
  Turing.MultiTapeTM.runFrom_succ_eq_step',
  Turing.moveInputPos_pos_of_ne_right, Turing.pairDiagCfg, Turing.pairDiagTM,
  Turing.pairDiag_inner, Turing.pairDiag_step}
\difficulty{4}
Let $x$ be a bit string, $w$ an output string, and $t \le |x|$.  Running the
six-state diagonal pairing machine for exactly $t$ steps from the
configuration in its copy state (state $5$) with input head at position $1$
and output $w$ yields the configuration still in the copy state with input
head at position $t + 1$ and output $w$ followed by the first $t$ bits of
$x$.
\end{lemma}

\begin{lemma}[Separator emission and turnaround]
\lean{Turing.pairDiag_separator}
\label{Turing.pairDiag_separator}
\leanok
\uses{Turing.Cfg.ext_zero_tapes, Turing.MultiTapeTM.runFrom,
  Turing.MultiTapeTM.runFrom_succ_eq_step', Turing.MultiTapeTM.runFrom_zero,
  Turing.moveInputPos_neg_of_ne_left, Turing.moveInputPos_zero,
  Turing.pairDiagCfg, Turing.pairDiagTM, Turing.pairDiag_right,
  Turing.pairDiag_step}
\difficulty{4}
Let $x$ be a bit string and $w$ an output string.  Running the six-state
diagonal pairing machine for exactly $3$ steps from the configuration in its
initial control state with input head at the right boundary position
$|x| + 1$ and output $w$ yields the configuration in its rewind state
(state $4$) with input head at position $|x|$ and output $w$ followed by the
separator pair $01$.
\end{lemma}

\begin{lemma}[The complete pairing run]
\lean{Turing.pairDiag_run}
\label{Turing.pairDiag_run}
\leanok
\uses{Turing.MultiTapeTM.initCfg, Turing.MultiTapeTM.runFrom,
  Turing.MultiTapeTM.runFrom_add, Turing.MultiTapeTM.runFrom_succ_eq_step',
  Turing.moveInputPos_zero, Turing.pairDiagCfg, Turing.pairDiagTM,
  Turing.pairDiag_copy, Turing.pairDiag_double, Turing.pairDiag_rewind,
  Turing.pairDiag_right, Turing.pairDiag_separator, Turing.pairDiag_step,
  Turing.pairEncode}
\difficulty{5}
For every bit string $x$ of length $n$, running the six-state diagonal
pairing machine from its initial configuration on input $x$ for exactly
$4n + 5$ steps yields a halted configuration with input head at position
$n + 1$ whose output is the diagonal pair encoding $\langle x, x\rangle$
(each bit of $x$ doubled, then the separator $01$, then $x$ verbatim).
\end{lemma}

\begin{theorem}[The diagonal pairing is linear-time computable]
\lean{Turing.computesFunInTime_pairEncode_diag}
\label{Turing.computesFunInTime_pairEncode_diag}
\leanok
\uses{Turing.FinTM, Turing.FinTM.ComputesFunInTime,
  Turing.FinTM.ComputesInTime, Turing.FinTM.ComputesInTime.mono,
  Turing.pairDiagTM, Turing.pairDiag_run, Turing.pairEncode}
\difficulty{3}
There exist a finite-state machine $M$ over the binary alphabet and a
constant $c$ such that on every input string $\alpha$ of length $n$, the
machine $M$ halts within $c \cdot (n + 1)$ steps with output the diagonal
pair encoding $\langle \alpha, \alpha\rangle$ (each bit of $\alpha$ doubled,
then the separator $01$, then $\alpha$ verbatim).  In particular the diagonal
pairing $\alpha \mapsto \langle \alpha, \alpha\rangle$ is computable in
linear time.
\end{theorem}

\begin{definition}[Serialization of a head move]
\lean{Turing.signBits}
\label{Turing.signBits}
\leanok
The fixed two-bit serialization of a head move: a left move is $11$, staying
put is $00$, and a right move is $10$.
\end{definition}

\begin{definition}[Serialization of an optional bit]
\lean{Turing.optBoolBits}
\label{Turing.optBoolBits}
\leanok
The fixed two-bit serialization of an optional bit: absent is $00$, the bit
$0$ is $10$, and the bit $1$ is $11$.
\end{definition}

\begin{definition}[Serialization of an optional write]
\lean{Turing.optOptBoolBits}
\label{Turing.optOptBoolBits}
\leanok
The fixed two-bit serialization of an optional write action on a binary
tape, where a present write may itself write the blank: no write is $00$,
writing the blank is $01$, writing the bit $0$ is $10$, and writing the bit
$1$ is $11$.
\end{definition}

\begin{definition}[Unary serialization of a state index]
\lean{Turing.unaryFin}
\label{Turing.unaryFin}
\leanok
The self-delimiting unary serialization of an index $s \in \{0,\dots,n-1\}$:
$s$ copies of the bit $1$ followed by a single bit $0$.
\end{definition}

\begin{definition}[Serialization of an optional successor state]
\lean{Turing.optStateBits}
\label{Turing.optStateBits}
\leanok
\uses{Turing.unaryFin}
The serialization of an optional successor state: halting (no successor) is
the single bit $0$, and a successor state $s$ is the bit $1$ followed by the
self-delimiting unary serialization of $s$.
\end{definition}

\begin{definition}[Serialization of a transition record]
\lean{Turing.actionBits}
\label{Turing.actionBits}
\leanok
\uses{Turing.Action, Turing.optBoolBits, Turing.optOptBoolBits,
  Turing.optStateBits, Turing.signBits}
The serialization of a single transition record of a one-work-tape binary
machine with $n + 1$ states: the concatenation of the two-bit input-head
move, the two-bit optional write on the work tape, the two-bit work-head
move, the two-bit optional output emission, and the self-delimiting
serialization of the optional successor state.
\end{definition}

\begin{definition}[Canonical serialization of a coded machine]
\lean{Turing.CodeTM.serialize}
\label{Turing.CodeTM.serialize}
\leanok
\uses{Turing.CodeTM, Turing.actionBits, Turing.pairEncode, Turing.unaryFin}
The fixed, scheme-independent canonical serialization of a machine in code
normal form: the pair encoding whose first component is the binary
representation of the state-count parameter (self-delimiting via the
doubled-bit region) and whose second component is the self-delimiting unary
serialization of the initial state followed by the serialized transition
records of the entire transition table, listed in a fixed enumeration order
(states in increasing order, with the scanned input symbol and the scanned
work symbol each ranging over blank, $0$, $1$).
\end{definition}

\begin{definition}[Machine representation scheme]
\lean{Turing.MachineCode}
\label{Turing.MachineCode}
\leanok
\uses{Turing.CodeTM}
The algebraic laws of a representation scheme for machines in code normal
form: an encoding map from coded machines to bit strings, a total decoding
map from bit strings to coded machines (so every string represents some
machine), and the law that the code of a machine followed by any number of
padding bits $1$ decodes back to that machine (so every machine has
infinitely many representations).
\end{definition}

\begin{theorem}[Decoding a code recovers the machine]
\lean{Turing.MachineCode.decode_encode}
\label{Turing.MachineCode.decode_encode}
\leanok
\uses{Turing.CodeTM, Turing.EffectiveMachineCode, Turing.MachineCode}
\difficulty{1}
For every representation scheme satisfying the algebraic laws and every
machine $M$ in code normal form, decoding the code of $M$ yields $M$ itself.
\end{theorem}

\begin{definition}[Effective representation scheme]
\lean{Turing.EffectiveMachineCode}
\label{Turing.EffectiveMachineCode}
\leanok
\uses{Turing.CodeTM.serialize, Turing.FinTM, Turing.FinTM.ComputesFunInTime,
  Turing.MachineCode}
An \emph{effective} representation scheme: a representation scheme
satisfying the algebraic laws, together with a \emph{canonizer} --- a
finite-state binary-alphabet machine --- and a time bound
$T : \mathbb{N} \to \mathbb{N}$ such that on every bit string $\alpha$ the
canonizer halts within $T(|\alpha|)$ steps and outputs the fixed canonical
serialization of the machine to which $\alpha$ decodes.
\end{definition}

\begin{theorem}[Every one-tape binary machine has a coded equivalent]
\lean{Turing.exists_codeTM}
\label{Turing.exists_codeTM}
\leanok
\uses{Turing.Cfg.mapState, Turing.CodeTM, Turing.CodeTM.toFinTM,
  Turing.FinTM, Turing.FinTM.ComputesInTime,
  Turing.MultiTapeTM.ComputesInTimeAndSpace, Turing.MultiTapeTM.relabelState,
  Turing.MultiTapeTM.relabelState_runFrom_init}
\difficulty{4}
Let $M$ be a machine over the binary alphabet with finitely many states and
exactly one work tape.  Then there exists a machine $M'$ in code normal form
such that for every input string $x$, every output string $w$, and every
time bound $t$, the bundled machine of $M'$ halts on $x$ within $t$ steps
with output $w$ if and only if $M$ does.
\end{theorem}
